\textbf{ЗМІСТ}

{\def\LTcaptype{none} % do not increment counter
\begin{longtable}[]{@{}
  >{\raggedright\arraybackslash}p{(\linewidth - 2\tabcolsep) * \real{0.9189}}
  >{\raggedleft\arraybackslash}p{(\linewidth - 2\tabcolsep) * \real{0.0811}}@{}}
\toprule\noalign{}
\endhead
\bottomrule\noalign{}
\endlastfoot
\textbf{АНОТАЦІЯ} & 2 \\
\textbf{ПЕРЕЛІК УМОВНИХ СКОРОЧЕНЬ} & 20 \\
\textbf{ВСТУП} & 21 \\
\textbf{РОЗДІЛ 1. ТЕОРЕТИЧНІ ЗАСАДИ НАВЧАННЯ РОЗРОБКИ ІМЕРСИВНИХ ОСВІТНІХ РЕСУРСІВ} & 33 \\
1.1. Розробка імерсивних освітніх ресурсів як педагогічна проблема підготовки майбутніх учителів інформатики. & 33 \\
1.2. Сутнісна характеристика імерсивних освітніх ресурсів. & 52 \\
1.3. Вітчизняний досвід використання імерсивних освітніх ресурсів

у професійній підготовці майбутніх учителів інформатики & 71 \\
\textbf{Висновки до першого розділу} & 84 \\
\textbf{РОЗДІЛ 2. ОРГАНІЗАЦІЙНО-ПЕДАГОГІЧНІ ОСНОВИ НАВЧАННЯ РОЗРОБКИ ІМЕРСИВНИХ ОСВІТНІХ РЕСУРСІВ} & 88 \\
2.1. Дидактичні підходи та принципи професійної підготовки майбутніх учителів інформатики до розробки імерсивних освітніх ресурсів & 86 \\
2.2. Структурно-функціональна модель підготовки майбутніх вчителів інформатики до розробки імерсивних освітніх ресурсів & 106 \\
2.3. Педагогічні умови навчання майбутніх вчителів інформатики розробки імерсивних освітніх ресурсів. & 117 \\
\textbf{Висновки до другого розділу} & 133 \\
\textbf{РОЗДІЛ 3. ОРГАНІЗАЦІЯ ТА МЕТОДИКА ПРОВЕДЕННЯ ДОСЛІДНО- ЕКСПЕРИМЕНТАЛЬНОЇ РОБОТИ} & 135 \\
3.1. Розробка критеріального апарату для оцінки \ldots.

Критерії, показники та рівні сформованості вмінь майбутніх учителів інформатики у сфері розробки імерсивних освітніх ресурсів. & 135 \\
3.2. Організація і методика проведення педагогічного експерименту & 150 \\
3.3. Аналіз результатів дослідно-експериментальної орботи

Результати та аналіз педагогічного експерименту & 162 \\
\textbf{Висновки до третього розділу} & 172 \\
\textbf{ВИСНОВКИ} & 175 \\
\textbf{СПИСОК ВИКОРИСТАНИХ ДЖЕРЕЛ} & 180 \\
\textbf{ДОДАТКИ} & 200 \\
\end{longtable}
}
